\section{Materials and Methods}
\label{sec:methodology}

\subsection{System Architecture and Operational Flow}
\label{sec:architecture}

The framework is built as a set of loosely coupled software services that generate candidate 4D trajectories and check them locally, for structured corridor operations within U-space. The services are part of a larger microservice fleet-control architecture for UAM operations~\cite{nagel2026microservice}. The main components for this paper are:
(i) a browser-based \emph{operator user interface} (UI), a web application (webapp) that the fleet operator uses to prepare and supervise flights,
(ii) a \emph{fleet webserver}, which works as backend-for-frontend (BFF) for the operator UI; it is the only service the browser talks to, keeps the local plan state, provides the vertiport and corridor network, and calls the other services on the server side,
(iii) a \emph{Planning Service}, which computes the route, synthesizes the waypoints, generates the 4D trajectory with CFMS, and re-plans a request when needed,
(iv) a \emph{Deconfliction Service}, which detects and reports conflicts with NDMap, and
(v) a \emph{USSP-facing operation-plan interface}, through which a locally prepared trajectory could later be submitted for U-space authorization, traffic information and lifecycle recording.

Two existing DLR tools are used inside these services. CFMS is an experimental flight-management and 4D trajectory generation system for unmanned multicopters, UAM aircraft and other VTOL aircraft. It takes a waypoint sequence with altitude, speed and flight-phase information and generates a dense, time-parameterized 4D trajectory that respects the flight dynamics of the aircraft~\cite{morscheck2021cfms}. NDMap is a conflict-detection and -resolution library for multi-dimensional objects. It stores trajectories in a combined space--time index and quickly finds where a new trajectory comes too close to a set of existing ones~\cite{NDMap}. In our framework, CFMS runs inside the Planning Service and NDMap runs inside the Deconfliction Service.

Trajectory generation and conflict checking are kept apart. The Planning Service creates the candidate trajectories and changes the graph for re-planning, while all NDMap-related work happens in the Deconfliction Service. This way the routing and trajectory code does not depend on the NDMap state, and each part can be deployed or scaled on its own.

The USSP interface follows the split of responsibilities used in U-space fleet-control architectures~\cite{nagel2026microservice}: the operator's planning system prepares the route, the candidate trajectory and a local conflict check, while the USSP stays responsible for U-space decisions such as plan authorization, traffic information, geo-awareness and lifecycle records. In the target workflow, the fleet webserver sends a versioned operation plan (U-plan) to the USSP. The USSP then runs its authorization workflow, which may also call the Deconfliction Service, and accepts the plan, asks for changes, or rejects it. The Planning and Deconfliction Services are only providers in this flow and do not authorize flights themselves~\cite{nagel2026microservice}. The USSP side of this interface is an assumption of the target concept; we did not implement or test it. Any change to the trajectory, time window, constraints, reservations or aircraft profile would need a new plan version. Our experiments stop before the authorization step, so the accepted trajectories in this paper are locally deconflicted candidate plans and not USSP-authorized flights.

\begin{figure}[!htbp]
  \centering
  \resizebox{0.92\textwidth}{!}{%
  \begin{tikzpicture}[
    box/.style={
      draw,
      rounded corners=8pt,
      align=center,
      minimum width=30mm,
      minimum height=14mm
    },
    req/.style={-Latex, thick},
    resp/.style={-Latex, thick, dashed},
    lbl/.style={font=\small, align=center},
    lbls/.style={font=\scriptsize, align=center, inner sep=1.5pt},
    node distance=26mm and 46mm
  ]

  % Nodes: operator side (top row), provider services (bottom row).
  % The two provider boxes are taller and further apart than the rest so the
  % four separate messages they exchange can each carry their own arrow.
  \node[box] (webapp) {Operator UI\\(webapp)};
  \node[box, right=of webapp] (bff) {Fleet BFF\\(webserver)};
  \node[box, right=of bff] (ussp) {U-space Service\\Provider (USSP)};

  \node[box, below=of bff, minimum height=32mm] (planner) {Planning Service\\(A* + CFMS)};
  \node[box, below=of ussp, minimum height=32mm] (deconf) {Deconfliction\\Service (NDMap)};

  % Operator UI talks only to the fleet BFF
  \draw[req] (webapp) -- node[lbl, above] {same-origin} node[lbl, below] {REST/WSS} (bff);

  % BFF submits versioned U-plans to the USSP
  \draw[req] (bff) -- node[lbl, above] {versioned U-plan\\submission} (ussp);

  % BFF requests route + trajectory from the Planning Service
  \draw[req] ([xshift=-4mm]bff.south) --
    node[lbl, left] {route intent,\\topology}
    ([xshift=-4mm]planner.north);
  \draw[resp] ([xshift=4mm]planner.north) --
    node[lbl, right] {candidate 4D\\trajectory}
    ([xshift=4mm]bff.south);

  % USSP consults the Deconfliction Service during authorization
  \draw[req] (ussp) --
    node[lbl, right] {deconfliction check\\at authorization\\(target concept)}
    (deconf);

  % Candidate conflict check evaluated in this paper. The two services exchange
  % four distinct messages; each direction gets its own arrow so that no arrow
  % stands for two operations travelling in opposite directions.
  \draw[req] ([yshift=13mm]planner.east) --
    node[lbls, above=0.8mm] {1. conflict check\\(candidate + peer ids)}
    ([yshift=13mm]deconf.west);
  \draw[req] ([yshift=4.5mm]deconf.west) --
    node[lbls, above=0.8mm] {2. trajectory fetch}
    ([yshift=4.5mm]planner.east);
  \draw[resp] ([yshift=-4.5mm]planner.east) --
    node[lbls, below=0.8mm] {3. trajectory samples}
    ([yshift=-4.5mm]deconf.west);
  \draw[resp] ([yshift=-13mm]deconf.west) --
    node[lbls, below=0.8mm] {4. conflict report}
    ([yshift=-13mm]planner.east);

  \end{tikzpicture}%
  }
  \caption{Service architecture and information flow. The operator UI in the browser talks only to the fleet webserver, which works as backend-for-frontend (BFF): it keeps the local plan state, provides the corridor network, asks the Planning Service for candidate 4D trajectories, and sends versioned U-plans to the USSP~\cite{nagel2026microservice}. REST (Representational State Transfer) is the request/response style of the service interfaces, and WSS (WebSocket Secure) is the encrypted two-way channel over which the fleet webserver sends plan updates to the operator UI. In the target concept, the USSP may call the Deconfliction Service during authorization; this interface is assumed and not evaluated. Solid arrows are requests, dashed arrows are responses, and each arrow points in the direction of the message. Between the Planning and Deconfliction Services, the trajectory fetch (2--3) is a separate exchange started by the Deconfliction Service, in the opposite direction to the conflict check (1) and its report (4). This paper evaluates this conflict check and the lateral re-routing that follows from it.}
  \label{fig:architecture}
\end{figure}

Figure~\ref{fig:architecture} shows how the services interact. Two points are important. First, the browser never calls the provider services directly; the fleet webserver is the only REST interface for the browser and calls the Planning Service, the Deconfliction Service and the USSP on the server side~\cite{nagel2026microservice}. Second, the conflict check studied in this paper runs between the Planning Service and the Deconfliction Service: the Planning Service sends a reference to the candidate trajectory together with the identifiers of the already accepted peer trajectories, and the Deconfliction Service then fetches these trajectories from the Planning Service's trajectory store before it runs the NDMap check. The internal planning loop, with conflict detection in the Deconfliction Service and lateral re-routing in the Planning Service, is shown in Figure~\ref{fig:planning_loop}.

\subsection{Global Route Planning on Corridor Graphs}
\label{sec:global_planning}

Routing between vertiports is a shortest-path problem on a weighted graph \(G=(V,E)\). The nodes \(V\) are vertiports and the edges \(E\) are corridor connections. The weight of an edge is the distance between its two end points, computed with the Haversine formula. We use A* search to find the shortest route between origin and destination.

As heuristic we use the straight-line distance from the current node to the destination. This heuristic never overestimates, because a route along graph edges cannot be shorter than the direct distance. The output of this stage is an ordered list of corridor nodes,
\[
P = [v_0, v_1, \dots, v_n],
\]
which is the lateral reference route for the next steps, waypoint synthesis and trajectory generation.

\subsection{4D Trajectory Generation via CFMS}
\label{sec:trajectory_generation}

The route \(P\) is turned into a four-dimensional trajectory (latitude, longitude, altitude, time) in two steps: waypoint synthesis and trajectory generation with CFMS~\cite{morscheck2021cfms}. Algorithm~\ref{alg:planner_pipeline} gives an overview of the planner pipeline; the steps are explained in the rest of this section.

\begin{Algorithm}[Planner pipeline for 4D trajectory generation]
\label{alg:planner_pipeline}
\begin{algorithmic}[1]
\REQUIRE Origin vertiport \(o\), destination vertiport \(d\), corridor graph \(G=(V,E)\), vertiport data, cruise altitude \(H_c\)
\ENSURE Dense trajectory \(\mathcal{T}_{\mathrm{dense}}\)
\STATE Compute corridor route \(P\) from \(o\) to \(d\) using A* on \(G\)
\STATE Generate phase-structured waypoint sequence \(W\) from \(P\), vertiport data, and \(H_c\)
\STATE Convert \(W\) into CFMS-compatible input
\STATE Generate dense 4D trajectory \(\mathcal{T}_{\mathrm{dense}}\) using CFMS
\RETURN \(\mathcal{T}_{\mathrm{dense}}\)
\end{algorithmic}
\end{Algorithm}

Figure~\ref{fig:waypoint-synthesis-schematic} shows the waypoint synthesis between A* routing and CFMS. A* gives the lateral sequence of corridor nodes, and the waypoint synthesizer adds terminal waypoints around the departure and arrival vertiports. The intermediate corridor nodes from A* are kept as cruise references and marked as navigation nodes in the waypoint list. The final waypoint sequence is then converted into CFMS input. In the figure, \(h_1\) and \(h_2\) are the lower and upper terminal reference heights, OFV is the obstacle-free volume around each vertiport, OLS the outer limitation surface, LDP the landing decision point, MSL mean sea level, and \(H_c\) the fixed cruise altitude. These terms are explained in more detail in the terminal-procedure part below.

\begin{figure}[!htbp]
\centering
\resizebox{\linewidth}{!}{%
\begin{tikzpicture}[
    >=Latex,
    font=\small,
    route/.style={very thick, blue!70!black},
    waypoint/.style={circle, draw=black, fill=white, inner sep=1.6pt},
    aux/.style={gray!55, dashed, thin},
    leader/.style={gray!65, thin},
    labelbox/.style={font=\scriptsize, align=center, fill=white, inner sep=1.8pt},
    smalllabel/.style={font=\scriptsize, align=center, fill=white, inner sep=1.5pt},
    volume/.style={fill=blue!10, draw=blue!35, dashed, rounded corners=2pt},
    axis/.style={gray!65, thin, -Latex}
]

% ------------------------------------------------------------------
% Coordinates
% ------------------------------------------------------------------
\coordinate (origin)   at (0,0);
\coordinate (depHone)  at (0,1.05);
\coordinate (depHtwo)  at (1.55,2.05);
\coordinate (ols)      at (2.65,3.15);
\coordinate (cruiseA)  at (3.55,4.20);
\coordinate (cruiseB)  at (5.80,4.20);
\coordinate (cruiseC)  at (8.05,4.20);
\coordinate (tod)      at (9.35,4.20);
\coordinate (arrEntry) at (10.65,3.05);
\coordinate (arrHtwo)  at (11.75,2.05);
\coordinate (arrHone)  at (11.75,1.05);
\coordinate (ldp)      at (11.75,0.55);
\coordinate (dest)     at (11.75,0);

% Useful intermediate targets
\coordinate (climbmid)      at ($(origin)!0.55!(depHone)$);
\coordinate (transitionmid) at ($(depHone)!0.62!(depHtwo)$);

% ------------------------------------------------------------------
% Background reference lines and axes
% ------------------------------------------------------------------
\draw[axis] (-0.85,-0.35) -- (13.10,-0.35);
\draw[axis] (-0.85,-0.35) -- (-0.85,5.05);

\node[rotate=90, gray!70, font=\scriptsize] at (-1.25,2.45) {altitude / MSL};
\node[gray!70, font=\scriptsize] at (6.15,-1.15)
{progress along selected A* corridor route};

\draw[aux] (-0.55,4.20) -- (12.75,4.20);
\node[labelbox, anchor=west] at (10.15,4.45) {$H_c = 200\,\mathrm{m}$ MSL};

\draw[aux] (-0.55,2.05) -- (0.60,2.05);
\draw[aux] (11.45,2.05) -- (12.95,2.05);
\node[smalllabel, font=\small, text=black, anchor=west] at (-0.48,2.22) {$h_{2}$};

\draw[aux] (-0.55,1.05) -- (0.45,1.05);
\draw[aux] (11.45,1.05) -- (12.95,1.05);
\node[smalllabel, font=\small, text=black, anchor=west] at (-0.48,1.22) {$h_{1}$};

\draw[aux] (-0.55,0.55) -- (0.45,0.55);
\draw[aux] (11.45,0.55) -- (12.95,0.55);

% ------------------------------------------------------------------
% Terminal OFV / verticone volumes
% ------------------------------------------------------------------
\begin{scope}[on background layer]
    \path[volume]
        (-0.18,0.00) --
        (0.62,0.00) --
        (1.05,1.45) --
        (1.80,2.55) --
        (1.35,2.80) --
        (-0.18,2.80) -- cycle;

    \path[volume]
        (11.15,0.00) --
        (12.30,0.00) --
        (12.30,2.80) --
        (10.75,2.80) --
        (11.25,2.10) --
        (11.55,1.10) -- cycle;
\end{scope}

% ------------------------------------------------------------------
% Main synthesized waypoint route
% ------------------------------------------------------------------
\draw[route]
    (origin) --
    (depHone) --
    (depHtwo) --
    (ols) --
    (cruiseA) --
    (cruiseB) --
    (cruiseC) --
    (tod) --
    (arrEntry) --
    (arrHtwo) --
    (arrHone) --
    (ldp) --
    (dest);

% Waypoint markers
\foreach \p in {origin,depHone,depHtwo,ols,cruiseA,cruiseB,cruiseC,tod,arrEntry,arrHtwo,arrHone,ldp,dest}
{
    \node[waypoint] at (\p) {};
}

% ------------------------------------------------------------------
% Labels: departure side
% ------------------------------------------------------------------
\node[labelbox, anchor=north] at (0,-0.48) {origin pad\\\(v_0\)};

\node[labelbox, anchor=east] (vertclimblabel) at (-0.45,0.72) {vertical\\climb};
\draw[leader] (vertclimblabel.east) -- (climbmid);

\node[labelbox] (depofvlabel) at (0.78,3.30) {departure OFV\\verticone};

\node[labelbox] (transitionlabel) at (2.35,1.48) {transition};
\draw[leader] (transitionlabel.west) to[out=180,in=15] (transitionmid);

\node[labelbox] (olslabel) at (2.35,4.95) {optional\\OLS point};
\draw[leader] (olslabel.south) to[out=-90,in=120] (ols);

% ------------------------------------------------------------------
% Labels: cruise / corridor middle
% ------------------------------------------------------------------
\node[labelbox] (cruiselabel) at (5.95,5.02) {intermediate corridor nodes\\cruise references};
\draw[leader] (cruiselabel.south west) to[out=-110,in=90] (cruiseB);
\draw[leader] (cruiselabel.south east) to[out=-70,in=90] (cruiseC);

\node[labelbox] (todlabel) at (9.35,5.02) {top of\\descent};
\draw[leader] (todlabel.south) -- (tod);

\node[
    draw=gray!55,
    rounded corners=2pt,
    fill=white,
    align=center,
    inner sep=4pt
] at (6.20,2.45)
{waypoint synthesis adds\\terminal and vertical structure\\[2pt]
CFMS-compatible waypoint sequence\\[-1pt]
\(\displaystyle W=[w_1,\ldots,w_m]\)};

% Height-reference legend: h1/h2 are pad-relative, H_c is absolute.
\node[
    draw=gray!55,
    rounded corners=2pt,
    fill=white,
    align=left,
    font=\scriptsize,
    inner sep=4pt
] at (6.20,0.72)
{\(h_{1}, h_{2}\), LDP, OLS: heights \emph{above the respective pad}\\
\(H_c\): absolute altitude above MSL};

\node[smalllabel] at (5.90,-0.58)
{\(v_1,\ldots,v_{n-1}\): intermediate A* corridor nodes of \(P=[v_0,\ldots,v_n]\)};

% Light gray guide dots for corridor-node sequence
\foreach \x in {1.7,3.55,5.80,8.05,9.35,10.65}
{
    \node[circle, fill=gray!45, inner sep=1.1pt] at (\x,-0.35) {};
}

% ------------------------------------------------------------------
% Labels: arrival side
% ------------------------------------------------------------------
\node[labelbox] (arrofvlabel) at (11.80,3.52) {arrival OFV\\verticone};

\node[labelbox, anchor=west] (ofventrylabel) at (12.55,2.38) {OFV entry};
\draw[leader] (ofventrylabel.west) to[out=180,in=0] (arrHtwo);

\node[labelbox, anchor=west] (overpadlabel) at (12.55,1.40) {over-pad\\waypoint};
\draw[leader] (overpadlabel.west) to[out=180,in=0] (arrHone);

\node[labelbox, anchor=west] (ldplabel) at (12.55,0.58) {LDP};
\draw[leader] (ldplabel.west) -- (ldp);

\node[labelbox, anchor=north] at (11.75,-0.48) {destination pad\\\(v_n\)};

\end{tikzpicture}%
}
\caption{Schematic waypoint-synthesis procedure used to convert an A*-selected corridor route into a CFMS-compatible phase-structured waypoint sequence. The route is \(P = [v_0, v_1, \dots, v_n]\) with origin \(v_0\) and destination \(v_n\); the intermediate corridor nodes \(v_1, \dots, v_{n-1}\) define cruise-route references at the fixed cruise layer. The departure and arrival procedures introduce terminal waypoints at \(h_1\), \(h_2\), and the LDP. The two height references differ: \(h_1\), \(h_2\), the LDP and the OLS target are heights above the respective departure or arrival pad, whereas \(H_c\) is an absolute altitude above MSL. Because pad elevations differ, equal \(h_1\) values at the two ends correspond to different absolute altitudes.}
\label{fig:waypoint-synthesis-schematic}
\end{figure}

\paragraph{Inputs.}
The trajectory generation stage consumes:
(i) the node sequence \(P\),
(ii) vertiport coordinates and elevations,
(iii) a fixed cruise altitude layer, and
(iv) nominal phase-dependent speeds.
In this study, the cruise layer is a fixed altitude of \SI{200}{m} above MSL. Because it is defined above MSL, the height above ground changes with terrain and pad elevation. We use this simplification on purpose, to test the planning stack under controlled conditions. The layer is higher than the height band usually assumed for normal low-altitude UAS operations. We chose it because UAM corridors are seen as a separate airspace structure, and because one fixed layer removes altitude assignment as an extra variable from the benchmark. We do not claim that \SI{200}{m} is the right corridor altitude for real operations. Terrain, obstacles and weather are not modeled in this step, so the trajectories are benchmark trajectories and not operational flight plans.

\paragraph{Nominal waypoint synthesis.}
The waypoint generator turns the lateral A* route \(P\) into a short waypoint sequence \(W=[w_1,\dots,w_m]\) with flight-phase labels for take-off, climb, cruise, descent and landing. The first and last waypoints are the origin and destination pads. In between, the synthesizer adds a few more waypoints that shape the vertical profile and give CFMS a route with clear flight phases.

For routes with several nodes, the intermediate corridor nodes are kept as cruise reference points. Their lateral position follows the graph path, and their altitude follows the nominal climb up to the cruise layer. For a direct connection between origin and destination, a smaller waypoint set is used, with one extra point on the connecting great circle so that departure and arrival connect smoothly.

At the cruise nodes, the A* graph nodes are used as centerline references in the waypoint list. The trajectory generation does not use a turn or swept-volume constraint based on the corridor width. So any smoothing, turning or speed change around these nodes comes from the trajectory generator itself and should not be read as a validated corridor-width turn model. Checking that the aircraft stays inside a lateral corridor envelope is outside the scope of this benchmark.

Each waypoint gets a flight phase and a target speed. The phase labels separate the terminal segments near the vertiports from the cruise and transition segments along the corridors. The nominal speeds are based on the performance limits of the aircraft, scaled down to keep a margin.

\paragraph{Terminal procedure construction.}
The departure and arrival parts of \(W\) come from a simplified model of the vertiport terminal procedures. At departure, the sequence starts at the pad, continues with a vertical climb, and then turns into the outbound corridor direction. Arrival works the same way in reverse: the aircraft is led from the inbound corridor segment to the destination pad along a constrained descent.

For these terminal segments, the waypoint generator uses a parameterized obstacle-free volume around each vertiport, which we call a ``verticone''. \emph{Verticone} is our own term for this construction and not an official term of the European Union Aviation Safety Agency (EASA). Simply put, the verticone is a protected three-dimensional volume above and around the pad: close to the pad the aircraft moves mostly vertically, and higher up the allowed lateral area gets wider. The geometry parameters are simplified and inspired by the EASA Prototype Technical Specifications for the Design of Visual Flight Rules (VFR) Vertiports (PTS-VPT-DSN), mainly the vertical take-off and landing procedure and the OFV definitions in Chapter~D, Subpart~2~\cite{easa2022_pts_vpt_dsn}. The terminal waypoints are meant to stay inside this simplified volume. We have not checked exact conformance with the EASA specification, and this is not a certification or an assessment of operational procedures.

At departure, the aircraft first climbs vertically from the pad to the lower terminal height \(h_1\). It then moves towards the second reference height \(h_2\) while turning into the outbound corridor direction. If there is enough horizontal distance, the climb continues along the same centerline up to the outer limitation surface at up to \SI{152}{m} above the pad, but not higher than the cruise altitude. Arrival is the reverse: the trajectory may first pass a top-of-descent point on the inbound centerline, then enters the OFV at \(h_2\), moves to an over-pad waypoint at \(h_1\), passes the LDP reference height, and ends on the pad.

The result is a small, well-structured set of nominal waypoints that keeps the lateral corridor route and adds the altitude changes needed before and after cruise. Contingency and non-nominal terminal procedures are not part of this evaluation.

\paragraph{CFMS execution.}
The waypoint list is converted into the CFMS input format with World Geodetic System 1984 (WGS-84) coordinates, altitudes, flight-phase labels and speeds. CFMS then generates a dense, time-parameterized trajectory \(\mathcal{T}_{\mathrm{dense}}\) that follows the given kinematic profile~\cite{morscheck2021cfms}. The trajectory-generation times in the evaluation are measured for the CFMS call and the dense trajectory generation together.

\subsection{Deconfliction Service Using NDMap}
\label{sec:conflict_detection}

Conflict detection runs in the separate Deconfliction Service. The service uses NDMap~\cite{NDMap} on the 4D state space \((\lambda,\varphi,H,t)\), i.e.\ longitude, latitude, altitude and time. NDMap indexes the referenced trajectories and finds the times at which the candidate trajectory is too close to one of the peer trajectories. This keeps all conflict-management state and NDMap data structures out of the Planning Service.

\paragraph{Conflict predicate.}
The trajectories enter the detector as discrete 4D samples
\[
s^{X}_{k} = \bigl(\mathbf{r}^{X}_{k},\, t^{X}_{k}\bigr),
\qquad
\mathbf{r}^{X}_{k} = \bigl(\lambda^{X}_{k}, \varphi^{X}_{k}, H^{X}_{k}\bigr),
\]
of trajectory \(X\). Here \(\mathbf{r}^{X}_{k}\) is the position at sample time \(t^{X}_{k}\), with longitude \(\lambda\) and latitude \(\varphi\) in WGS-84 degrees (converted to radians inside the detector), altitude \(H\) in meters MSL, and \(t\) in epoch seconds. CFMS produces a sample about every \SI{1}{s}. The detector does not compare a sample of one trajectory directly with a sample of the other. Instead, it evaluates both trajectories at the same time \(t\) and gets each position by linear interpolation between the two samples around \(t\),
\begin{equation}
\label{eq:traj_interpolation}
\mathbf{r}^{X}(t) = (1-\alpha)\,\mathbf{r}^{X}_{k} + \alpha\,\mathbf{r}^{X}_{k+1},
\qquad
\alpha = \frac{t - t^{X}_{k}}{t^{X}_{k+1} - t^{X}_{k}},
\qquad
t \in \bigl[t^{X}_{k},\, t^{X}_{k+1}\bigr],
\end{equation}
which gives the interpolated position \(\mathbf{r}^{X}(t) = (\lambda^{X}(t), \varphi^{X}(t), H^{X}(t))\). Two trajectories \(A\) and \(B\) are in conflict at time \(t\) if and only if
\begin{equation}
\label{eq:conflict_predicate}
\begin{aligned}
\mathcal{C}_{AB}(t) \iff\;
  & t \in \mathcal{T}_{AB} \\
  &\wedge\; d_{\mathrm{geo}}\bigl((\lambda^{A}(t), \varphi^{A}(t)),\,
      (\lambda^{B}(t), \varphi^{B}(t))\bigr) < d_{\mathrm{hor}} \\
  &\wedge\; \bigl\lvert H^{A}(t) - H^{B}(t) \bigr\rvert \le d_{\mathrm{vert}},
\end{aligned}
\end{equation}
where \(\mathcal{T}_{AB} = [t^{A}_{0}, t^{A}_{\mathrm{end}}] \cap [t^{B}_{0}, t^{B}_{\mathrm{end}}]\) is the time period in which both trajectories exist. Only longitude and latitude go into the horizontal distance \(d_{\mathrm{geo}}\); altitude is compared separately. \(d_{\mathrm{geo}}\) is the distance on the WGS-84 ellipsoid, computed with Vincenty's inverse formula. The thresholds are \(d_{\mathrm{hor}} = 0.3\,\mathrm{NM} \approx \SI{555.6}{m}\), with NM meaning nautical miles, and \(d_{\mathrm{vert}} = 100\,\mathrm{ft} = \SI{30.48}{m}\). The horizontal comparison is strict and the vertical one is not, as written in Equation~(\ref{eq:conflict_predicate}), and unit tests check both thresholds from both sides. The vertical threshold is set in feet and converted to meters once when the detector is configured, so the comparison itself works only in meters, like the trajectory altitudes. Note that the route search of Section~\ref{sec:global_planning} uses a different and simpler distance, the spherical Haversine formula, which has nothing to do with Equation~(\ref{eq:conflict_predicate}).

Both thresholds are a configuration choice for this benchmark and not a proposal for UAM separation minima. They follow the usual air-traffic practice of giving minima in nautical miles and hundreds of feet, and we use them only to give the detector a fixed and repeatable criterion. Which separation minima are right for eVTOL aircraft in low-altitude corridors is still an open question and not part of this paper. Since both values are detector parameters, the experiments can be repeated with other values.

The fixed cruise layer of Section~\ref{sec:trajectory_generation} has one important effect. Since all trajectories cruise at the same altitude, two aircraft in cruise always meet the vertical criterion of Equation~(\ref{eq:conflict_predicate}). So in our benchmark, en-route conflicts are decided only by horizontal distance, and the vertical criterion only matters where the aircraft are at different heights, i.e.\ during climb and descent near the vertiports. This comes from the single-layer benchmark and not from the detector, which always checks both criteria.

Because both trajectories are interpolated to the same time before the check, there is no time difference between the two compared positions, and no time tolerance is needed.

\paragraph{From the predicate to a reported conflict.}
Two more steps happen before and after Equation~(\ref{eq:conflict_predicate}). Before, a pair of trajectories is only looked at during \(\mathcal{T}_{AB}\): if the two time periods do not overlap, the pair is skipped. Overlap in time is therefore only a precondition for checking a pair and not a conflict by itself---two trajectories can share their whole time period and still be conflict-free. After, the times at which the condition holds are grouped into conflict intervals \([t_{\text{start}}, t_{\text{end}}]\); violations that are \SI{2}{s} or less apart are merged into one interval. For each interval, the Deconfliction Service finds the closest approach, i.e.\ the time of smallest horizontal distance within the interval, and reports this time together with the horizontal and vertical distance at that moment. The interval and its closest-approach time are the input for the lateral re-routing in Section~\ref{sec:conflict_resolution}.

\paragraph{Sampling and completeness.}
At a cruise speed of about \SI{25}{m/s}, two samples are about \SI{25}{m} apart, which is small compared with the \SI{555.6}{m} horizontal threshold, and linear interpolation fits the straight, constant-speed cruise segments well. Still, the detector does not give a formal guarantee that it finds every conflict in continuous time, because it checks sampled times and does not report violations shorter than \SI{1}{s}.

We can estimate the effect for a simple case. Take two aircraft moving in straight lines with a constant horizontal relative speed \(v_{\mathrm{rel}}\) and a minimum horizontal distance \(d_{\min} < d_{\mathrm{hor}}\), and assume that the vertical criterion holds the whole time. Then the horizontal criterion holds for \(2\sqrt{d_{\mathrm{hor}}^{2} - d_{\min}^{2}}/v_{\mathrm{rel}}\), which goes to zero when \(d_{\min}\) gets close to \(d_{\mathrm{hor}}\). At \(v_{\mathrm{rel}} = \SI{50}{m/s}\), an encounter with \(d_{\min} \gtrsim \SI{555.0}{m}\) lasts less than \SI{1}{s} and is not reported. Real encounters have turns, speed changes and vertical motion, so this is only a rough estimate for encounters that just touch the threshold. These missed cases are the ones closest to the threshold and not deep violations, but the detector should be understood as reporting sampled violations of at least \SI{1}{s}, not as a proof of separation in continuous time. The unit tests cover this, including a grazing encounter below the \SI{1}{s} limit.

\paragraph{Detector state lifecycle.}
In the current version, the detector state is built again for every conflict check and is not kept between requests. Every check fetches the candidate and all peer trajectories from the Planning Service's trajectory store over the Hypertext Transfer Protocol (HTTP), builds the detector state, and throws it away after the check. The Deconfliction Service therefore keeps no state between requests. The candidate keeps the same trajectory identifier over all re-planning attempts, and each attempt overwrites the stored trajectory before the next check, so an old candidate cannot show up in later checks. Duplicates and conflicts of a trajectory with itself are avoided by removing duplicate identifiers and by only reporting conflicts that involve the candidate. The downside of this design is that each check gets more expensive as the peer set grows; the upside is that no state can get out of sync, even with several workers. Beyond a single host, a shared trajectory store would be needed. The effect on latency is shown in Section~\ref{sec:conflict_resolution_results}.

\begin{figure}[!htbp]
\centering
\begin{tikzpicture}[
  node distance=8mm,
  box/.style={
    draw,
    rounded corners,
    align=center,
    inner sep=4pt,
    minimum width=42mm,
    minimum height=8mm
  },
  arrow/.style={-Latex, thick}
]

\node[box] (fetch) {Route request from the fleet BFF\\(topology, peer-trajectory references)};
\node[box, below=of fetch] (astar) {Planning Service:\\Global routing (A*)};
\node[box, below=of astar] (traj) {Planning Service:\\4D trajectory generation (CFMS)};
\node[box, below=of traj] (deconf) {Deconfliction Service:\\Conflict check (NDMap)};
\node[box, below=of deconf] (submit) {Candidate plan for U-plan\\submission by the fleet BFF};

\node[box, right=20mm of deconf] (penalize) {Planning Service:\\Map conflict to edge\\and re-plan};

% Main vertical flow
\draw[arrow] (fetch) -- (astar);
\draw[arrow] (astar) -- (traj);
\draw[arrow] (traj) -- node[right, font=\small, align=center] {trajectory reference\\+ peer identifiers} (deconf);
\draw[arrow] (deconf) -- node[right, font=\small] {no conflict} (submit);

% Conflict loop
\draw[arrow] (deconf.east) -- node[above, font=\small] {conflict} (penalize.west);
\draw[arrow] (penalize.north) |- (astar.east);

\end{tikzpicture}
\caption{Operational planning loop up to the USSP-facing authorization boundary. The fleet BFF (webserver) issues the route request; candidate trajectories are generated by the Planning Service and checked by the Deconfliction Service, which fetches the referenced candidate and peer trajectories from the Planning Service's trajectory store before the NDMap check. Conflict reports trigger request-local edge removal, lateral re-routing, and re-execution of A* inside the Planning Service. A conflict-free candidate is returned to the fleet BFF for versioned U-plan submission to the USSP.}
\label{fig:planning_loop}
\end{figure}

\subsection{Conflict Resolution by Lateral Re-routing}
\label{sec:conflict_resolution}

Conflicts are resolved by \emph{lateral re-routing} on the corridor graph. The route segments that belong to a detected conflict are blocked for the current request, and A* is run again to find another route. The steps are:
(1) the Planning Service generates a candidate route using A*,
(2) the corresponding 4D trajectory is generated and checked by the Deconfliction Service using NDMap,
(3) if there are conflicts, the reported conflict times are mapped back to positions on the candidate trajectory,
(4) for each of these positions, the nearest edge of the current A* route is found,
(5) these directed edges are removed from the request's own copy of the graph by giving them an infinite weight,
\[
  d(u,v) := \infty, \qquad (u,v)\in E_{\text{conflict}},
\]
and
(6) A* is re-executed on the modified graph.
The loop stops when NDMap reports no conflict, when no new conflicting edge is found, when no route is left, or when the maximum number of attempts is reached. Figure~\ref{fig:conflict_resolution_re-routing} shows steps (4)--(6) on a small example graph.

\paragraph{Conflict-to-edge mapping.}
Steps (3)--(5) work as follows. For each conflict, the mapping time \(t^{*}\) is the reported closest-approach time from Section~\ref{sec:conflict_detection}. If it is missing, the middle of the interval \((t_{\text{start}} + t_{\text{end}})/2\) is used, and otherwise \(t_{\text{start}}\); a conflict without any time is skipped and counted as a mapping failure. For conflicts between two trajectories, the closest-approach time is always there, so the mapped point is where the candidate is closest to the other aircraft. We take the candidate sample closest in time to \(t^{*}\) and compare its position with every directed edge \((u,v)\) of the current route, using the point-to-segment distance in a local equirectangular projection. The closest edge is chosen; if a sample lies exactly on a shared node and two edges are equally close, the earlier edge in the route is kept. Only the flown direction \((u,v)\) is removed; the opposite direction \((v,u)\) stays available. If several conflicts map to the same edge, it is blocked only once. There is no maximum mapping distance. Because the mapped sample lies on the candidate trajectory, and the route is the lateral reference of that trajectory, the distances stay small in practice, limited by the terminal procedure geometry, but we do not enforce a limit, which is a known limitation. If an attempt produces no edge that is not already blocked, the loop ends with a separate \emph{mapping-stall} outcome (\texttt{stalled\_no\_new\_edges}), which we report separately from the retry-limit and no-route outcomes. Unit tests cover the tie-breaking, duplicate conflicts, edge direction and mapping failures.

\begin{figure}[!htbp]
    \centering
    \subfloat[Initial A* candidate path \(s \to u \to v \to d\).]{
    \begin{minipage}[t]{0.46\textwidth}
        \centering
        \begin{tikzpicture}[
            node distance=1.75cm, thick, >=Latex,
            avail/.style={-{Latex[length=2mm]}, gray!65, semithick},
            sel/.style={-{Latex[length=3mm]}, green!55!black, ultra thick},
            alt/.style={-{Latex[length=3mm]}, blue, ultra thick, dashed},
            blocked/.style={-{Latex[length=3mm]}, red, ultra thick}
        ]
            \node[circle,draw] (s) {\(s\)};
            \node[circle,draw,right of=s] (a) {\(u\)};
            \node[circle,draw,right of=a] (b) {\(v\)};
            \node[circle,draw,right of=b] (d) {\(d\)};
            \node[circle,draw,above of=b] (c) {\(w\)};

            % Every corridor connection carries a directed edge in each direction.
            \draw[avail] ([yshift=1.1mm]a.east) -- ([yshift=1.1mm]b.west);
            \draw[avail] ([yshift=-1.1mm]b.west) -- ([yshift=-1.1mm]a.east);
            \draw[avail] (s) -- (a);
            \draw[avail] (b) -- (d);
            \draw[avail] (a) -- (c);
            \draw[avail] (c) -- (d);

            \draw[sel] (s) -- (a);
            \draw[sel] ([yshift=1.1mm]a.east) -- ([yshift=1.1mm]b.west);
            \draw[sel] (b) -- (d);
        \end{tikzpicture}
    \end{minipage}}
    \hfill
    \subfloat[Only the flown direction \((u,v)\) is blocked; \((v,u)\) remains available and A* re-routes via \(w\).]{
    \begin{minipage}[t]{0.46\textwidth}
        \centering
        \begin{tikzpicture}[
            node distance=1.75cm, thick, >=Latex,
            avail/.style={-{Latex[length=2mm]}, gray!65, semithick},
            alt/.style={-{Latex[length=3mm]}, blue, ultra thick, dashed},
            blocked/.style={-{Latex[length=3mm]}, red, ultra thick}
        ]
            \node[circle,draw] (s) {\(s\)};
            \node[circle,draw,right of=s] (a) {\(u\)};
            \node[circle,draw,right of=a] (b) {\(v\)};
            \node[circle,draw,right of=b] (d) {\(d\)};
            \node[circle,draw,above of=b] (c) {\(w\)};

            \draw[avail] (s) -- (a);
            \draw[avail] (b) -- (d);
            \draw[avail] (a) -- (c);
            \draw[avail] (c) -- (d);
            % Reverse direction stays usable.
            \draw[avail] ([yshift=-1.1mm]b.west) -- ([yshift=-1.1mm]a.east);
            % Forward direction removed for this request only.
            \draw[blocked] ([yshift=1.1mm]a.east) -- ([yshift=1.1mm]b.west);
            \node[red, font=\footnotesize] at ($(a)!0.5!(b)+(0,0.75)$) {\(\times\)};

            \draw[alt] (s) -- (a);
            \draw[alt] (a) -- (c);
            \draw[alt] (c) -- (d);
        \end{tikzpicture}
    \end{minipage}}

    \vspace{2mm}
    {\footnotesize
    \textcolor{gray!65}{\textbf{---\!\!\(\rightarrow\)}} available directed edge \quad
    \textcolor{green!55!black}{\textbf{---\!\!\(\rightarrow\)}} initially selected route \quad
    \textcolor{red}{\textbf{---\!\!\(\rightarrow\)}} blocked direction \quad
    \textcolor{blue}{\textbf{-\,-\!\(\rightarrow\)}} re-routed path}

    \caption{Illustration of lateral re-routing for conflict resolution. Each undirected corridor connection is represented by one directed edge per travel direction. A conflict detected on the initial candidate trajectory is mapped to the directed edge \((u,v)\) it was flown on; only that direction is given effectively infinite weight for the current request, while the opposite direction \((v,u)\) and the persistent topology are left unchanged. A* then searches for an alternative path.}
    \label{fig:conflict_resolution_re-routing}
\end{figure}

NDMap can also support resolution in time, direction and altitude~\cite{NDMap}. In our framework, however, NDMap is only used for detection, and the Planning Service changes a request-local copy of the corridor graph to find another lateral route. The current resolution uses only lateral re-routing; it does not shift the departure time, change the speed, or move the aircraft to another altitude layer.

We kept this strategy simple on purpose. It never changes the persistent network: removed edges only apply to the request that is being re-planned. The downside is that re-routed paths can be longer, and some conflicts cannot be solved laterally at all, when all useful alternatives are blocked or the graph falls apart. Resolution in time or altitude is left for future work.

The next section describes the benchmark scenario, the system configuration and how we measured the timings.
